\begin{frame}{3분할이 실제로 깨지는 흔한 실수 세 가지}
  \begin{enumerate}
    \item \kb{테스트로 여러 모델을 비교}해 가장 잘 나온 걸 고르는 것
      (``살짝 훔쳐보기'').
    \item \kb{검증 데이터 없이} 테스트 데이터로 하이퍼파라미터까지 같이
      튜닝(검증과 테스트를 사실상 섞어버림).
    \item 데이터를 나누기 \textbf{전에} 정규화(스케일링) 같은 전처리
      통계량(평균, 표준편차)을 \textbf{전체 데이터}로 계산 — test 정보가
      전처리 단계에서 이미 학습에 스며듦(\kb{data leakage}).
  \end{enumerate}
  \vskip0.4em
  세 실수가 모두 같은 범주 — \textbf{``test 정보가 다른 곳으로 흘러나감''} —
  에 속한다. 겉모양은 달라(모델 고르기 / 전처리 순서 / …) 보여도
  \textbf{다 같은 정보 흐름의 역주행}.
  \vskip0.3em
  \begin{block}{더 좋은 질문}
    ``이 코드에 leakage가 있는가?''보다:
    \kq{``test(또는 val)의 정보가 학습이나 선택 단계에 닿았는가?''}
    \vskip0.2em
    아래 §에서 이 질문을 \textbf{숫자로} 답할 수 있는 도구를 하나씩 만든다.
  \end{block}
  {\footnotesize Week 8, 16 팀 프로젝트와 ML2에서도 이 3분할 원칙을 그대로
  지켜야 한다.}
\end{frame}
